\documentclass{article}
\usepackage{amsmath}
\usepackage{amssymb}
\usepackage{geometry}
\usepackage{graphicx}
\usepackage{xepersian}
\settextfont{XB Niloofar}

\geometry{a4paper, margin=1in}

\title{\lr{Bayesian Decision Theory} \\ (تئوری تصمیم‌گیری بیز)}
\author{}
\date{}
\begin{figure}[h!]
  \centering
  \includegraphics[width=0.75\textwidth]{img_page_15_1.png}
  \caption{توزیع‌های نرمال گوسی با ماتریس‌های کوواریانس برابر و قطری}
\end{figure}

\begin{figure}[h!]
  \centering
  \includegraphics[width=0.75\textwidth]{img_page_20_1.png}
  \caption{تغییر مرز تصمیم در اثر تغییر نسبت احتمالات پیشین کلاس‌ها}
\end{figure}






\begin{document}

\maketitle

\section*{فهرست مطالب}
\begin{itemize}
    \item مقدمه
    \item مفاهیم کلیدی
    \item فرمول بیز
    \item قاعده تصمیم‌گیری
    \item تعمیم به بیش از دو کلاس
    \item \lr{Discriminant Functions} (توابع تشخیص‌دهنده)
    \item \lr{Normal Density} (چگالی نرمال)
    \item \lr{Multivariate Normal Density} (چگالی نرمال چندمتغیره)
    \item تعاریف رسمی
    \item ویژگی‌های \lr{Covariance Matrix} (ماتریس کوواریانس)
    \item \lr{Discriminant Functions} برای \lr{Normal Density}
    \item خلاصه
\end{itemize}

\newpage

\section*{مقدمه}
\begin{itemize}
    \item \lr{Bayesian Decision Theory} یک رویکرد آماری برای \lr{pattern classification} (دسته‌بندی الگو) است.
    \item این تئوری موازنه بین تصمیمات دسته‌بندی را با استفاده از احتمالات و هزینه‌ها کمّی می‌کند.
    \item فرض می‌کند تمام احتمالات مربوطه شناخته شده هستند.
\end{itemize}
\begin{figure}[h!]
  \centering
  \includegraphics[width=0.75\textwidth]{img_page_2_1.jpeg}
  \caption{نمودار جریان تصمیم‌گیری بیزی و هزینه طبقه‌بندی}
\end{figure}

\begin{figure}[h!]
  \centering
  \includegraphics[width=0.75\textwidth]{img_page_4_1.png}
  \caption{توابع چگالی احتمال کلاس‌شرطی برای دو کلاس ماهی خاویار و سالمون}
\end{figure}

\begin{figure}[h!]
  \centering
  \includegraphics[width=0.75\textwidth]{img_page_6_1.png}
  \caption{تعمیم به طبقه‌بندی با بیش از دو کلاس و توابع متمایزکننده}
\end{figure}








مثال: دسته‌بندی ماهی (ماهی خاردار در مقابل ماهی سالمون) بر اساس ویژگی‌هایی مانند روشنایی.

\section*{مفاهیم کلیدی}
\begin{itemize}
    \item \lr{State of Nature} ($\omega$): نشان‌دهنده طبقه واقعی است (به عنوان مثال، ماهی خاردار یا ماهی سالمون).
    \item \lr{Prior Probability} ($P(\omega_j)$): احتمال یک طبقه قبل از مشاهده هرگونه داده.
    \item \lr{Class-Conditional Probability Density} ($p(x|\omega_j)$): احتمال مشاهده ویژگی $x$ به شرط طبقه $\omega_j$.
    \item \lr{Posterior Probability} ($P(\omega_j|x)$): احتمال طبقه $\omega_j$ پس از مشاهده ویژگی $x$.
\end{itemize}

\section*{فرمول بیز}
قضیه بیز (\lr{Bayes' Theorem}):
\[
\text{\lr{Posterior}} = \frac{\text{\lr{Likelihood}} \times \text{\lr{Prior}}}{\text{\lr{Evidence}}}
\]
یا به صورت ریاضی:
\[
P(\omega_j|x) = \frac{p(x|\omega_j)P(\omega_j)}{p(x)}
\]
که در آن \lr{Evidence} ($p(x)$) یک ثابت نرمال‌سازی است که اغلب در دسته‌بندی نادیده گرفته می‌شود و به صورت زیر محاسبه می‌شود:
\[
p(x) = \sum_{j=1}^{c} p(x|\omega_j)P(\omega_j)
\]

احتمال $P(\omega_j)$ به احتمال \lr{a posteriori} (یا احتمال \lr{posterior}) $P(\omega_j|x)$ تبدیل می‌شود — یعنی احتمال اینکه حالت طبیعت $\omega_j$ باشد، با فرض اینکه مقدار ویژگی $x$ اندازه‌گیری شده است. ما $p(x|\omega_j)$ را \lr{likelihood} مربوط به $\omega_j$ نسبت به $x$ می‌نامیم (اصطلاحی که برای نشان دادن این انتخاب شده است که در شرایط یکسان، طبقه‌ای که برای آن $p(x|\omega_j)$ بزرگ است، احتمال بیشتری دارد که طبقه واقعی باشد).

توجه داشته باشید که حاصل‌ضرب \lr{likelihood} و \lr{prior probability} است که در تعیین \lr{posterior probability} بیشترین اهمیت را دارد؛ فاکتور \lr{evidence} یعنی $p(x)$ را می‌توان صرفاً به عنوان یک فاکتور مقیاس‌دهی در نظر گرفت که تضمین می‌کند مجموع احتمالات پسین برابر با یک می‌شود، همانطور که تمام احتمالات خوب باید باشند. تغییرات $P(\omega_j|x)$ با $x$ در شکل ۲.۲ برای حالت $P(\omega_1) = \frac{2}{3}$ و $P(\omega_2) = \frac{1}{3}$ نشان داده شده است.

شکل ۲.۱: توابع چگالی احتمال شرطی کلاس فرضی، چگالی احتمال اندازه‌گیری یک مقدار ویژگی خاص $x$ را به شرط قرار داشتن الگو در طبقه $\omega_i$ نشان می‌دهند. اگر $x$ نشان‌دهنده طول ماهی باشد، این دو منحنی ممکن است تفاوت طول جمعیت دو نوع ماهی را توصیف کنند. توابع چگالی نرمال‌سازی شده‌اند و بنابراین مساحت زیر هر منحنی ۱.۰ است.

\section*{قاعده تصمیم‌گیری}
دسته‌بندی با کمترین نرخ خطا (\lr{Minimum Error Rate Classification}):
\begin{itemize}
    \item تصمیم بگیر $\omega_1$ اگر $P(\omega_1|x) > P(\omega_2|x)$ است؛ در غیر این صورت تصمیم بگیر $\omega_2$.
    \item معادل است با: تصمیم بگیر $\omega_1$ اگر $p(x|\omega_1)P(\omega_1) > p(x|\omega_2)P(\omega_2)$ است.
\end{itemize}
احتمال خطا به صورت زیر تعریف می‌شود:
\[
P(\text{error}|x) = \min[P(\omega_1|x), P(\omega_2|x)]
\]

شکل ۲.۲: احتمالات پسین برای پیشین‌های خاص $P(\omega_1) = \frac{2}{3}$ و $P(\omega_2) = \frac{1}{3}$ برای چگالی‌های احتمال شرطی کلاس نشان داده شده در شکل ۲.۱. بنابراین، در این حالت، با فرض اینکه یک الگو دارای مقدار ویژگی $x=14$ اندازه‌گیری شود، احتمال قرار داشتن آن در طبقه $\omega_2$ تقریباً $0.08$ و احتمال قرار داشتن آن در $\omega_1$ تقریباً $0.92$ است. در هر $x$، مجموع احتمالات پسین برابر با $1.0$ است.

برخی بینش‌های اضافی را می‌توان با بررسی چند حالت خاص به دست آورد:
\begin{enumerate}
    \item \textbf{حالت ۱:} اگر برای برخی از $x$ها داشته باشیم $p(x|\omega_1) = p(x|\omega_2)$، در این صورت آن مشاهده خاص هیچ اطلاعاتی درباره حالت طبیعت به ما نمی‌دهد. در این حالت، تصمیم‌گیری کاملاً به احتمالات پیشین $P(\omega_1)$ و $P(\omega_2)$ بستگی دارد.
    \item \textbf{حالت ۲:} اگر $P(\omega_1) = P(\omega_2)$ باشد، در این صورت حالت‌های طبیعت به یک اندازه محتمل هستند. در این حالت، تصمیم‌گیری کاملاً بر اساس درست‌نمایی‌ها (\lr{likelihoods}) $p(x|\omega_j)$ انجام می‌شود.
    \item \textbf{حالت کلی:} به طور کلی، هم احتمالات پیشین و هم درست‌نمایی‌ها در اتخاذ تصمیم مهم هستند. قاعده تصمیم‌گیری بیز (\lr{Bayes decision rule}) این فاکتورها را ترکیب می‌کند تا به حداقل احتمال خطا دست یابد.
\end{enumerate}

\section*{تعمیم به بیش از دو کلاس}
قاعده تصمیم‌گیری بیز برای کمینه‌سازی ریسک، خواستار انتخاب عملی است که ریسک شرطی را کمینه کند. برای کمینه‌سازی میانگین احتمال خطا، باید کلاسی ($\omega_i$) را انتخاب کنیم که احتمال پسین $P(\omega_i|x)$ را بیشینه می‌کند. به عبارت دیگر، برای حداقل نرخ خطا:
\[
\text{تصمیم بگیر } \omega_i \text{ اگر } P(\omega_i|x) > P(\omega_j|x) \text{ برای همه } j \neq i.
\]

این قاعده به طور طبیعی به کلاس‌های متعدد ($c > 2$) تعمیم می‌یابد. برای هر کلاس $\omega_i$، احتمال پسین $P(\omega_i|x)$ را محاسبه کرده و بردار ویژگی $x$ را به کلاسی با بیشترین احتمال پسین اختصاص می‌دهیم. این کار تضمین می‌کند که احتمال خطا کمینه شود.

\section*{توابع تشخیص‌دهنده (\lr{Discriminant Functions})}
تابع تشخیص‌دهنده ($g_i(x)$) برای اختصاص دادن یک بردار ویژگی $x$ به کلاس $\omega_i$ استفاده می‌شود.
برای حداقل نرخ خطا:
\[
g_i(x) = P(\omega_i|x)
\]
همچنین می‌تواند به صورت زیر بیان شود:
\[
g_i(x) = p(x|\omega_i)P(\omega_i)
\]
یا به فرم لگاریتمی:
\[
g_i(x) = \ln p(x|\omega_i) + \ln P(\omega_i)
\]

شکل ۲.۵: ساختار عملکردی یک دسته‌بندی‌کننده الگوی آماری عمومی که شامل ورودی‌ها و $c$ تابع تشخیص‌دهنده $g_i(x)$ است. در مرحله بعد مشخص می‌شود که کدام یک از مقادیر تشخیص‌دهنده حداکثر است و الگوی ورودی را بر این اساس دسته‌بندی می‌کند. فلش‌ها جهت جریان اطلاعات را نشان می‌دهند، اگرچه اغلب زمانی که جهت جریان واضح است، فلش‌ها حذف می‌شوند.

شکل ۲.۶: در این دسته‌بندی‌کننده دو طبقه دو بعدی، چگالی‌های احتمال گاوسی هستند (با بیضی‌های $1/e$ نشان داده شده‌اند)، مرز تصمیم‌گیری شامل دو هذلولی است و بنابراین ناحیه تصمیم‌گیری $R_1$ به سادگی همبند نیست.

\section*{چگالی نرمال (\lr{Normal Density})}
چگالی نرمال تک‌متغیره (\lr{Univariate Normal Density}) به صورت زیر است:
\[
p(x) = \frac{1}{\sqrt{2\pi}\sigma} \exp\left[ -\frac{1}{2}\left(\frac{x-\mu}{\sigma}\right)^2 \right]
\]
که در آن مقدار مورد انتظار $x$ (یک میانگین، در اینجا روی فضای ویژگی گرفته شده است) برابر است با:
\[
\mu = E[x] = \int_{-\infty}^{\infty} x p(x) dx
\]
و انحراف مربع مورد انتظار یا واریانس برابر است با:
\[
\sigma^2 = E[(x-\mu)^2] = \int_{-\infty}^{\infty} (x-\mu)^2 p(x) dx
\]
\begin{figure}[h!]
  \centering
  \includegraphics[width=0.75\textwidth]{img_page_7_1.png}
  \caption{ساختار عمومی یک سیستم طبقه‌بند آماری الگو}
\end{figure}

\begin{figure}[h!]
  \centering
  \includegraphics[width=0.75\textwidth]{img_page_8_1.png}
  \caption{توزیع نرمال یک‌متغیره و انحراف معیارهای آن}
\end{figure}






چگالی نرمال تک‌متغیره به طور کامل با دو پارامتر مشخص می‌شود: میانگین $\mu$ و واریانس $\sigma^2$. برای سادگی، اغلب $p(x)$ را با نوشتن $p(x) \sim N(\mu, \sigma^2)$ خلاصه می‌کنیم تا بگوییم $x$ به طور نرمال با میانگین $\mu$ و واریانس $\sigma^2$ توزیع شده است. نمونه‌های حاصل از توزیع‌های نرمال تمایل دارند در اطراف میانگین خوشه شوند، با پراکندگی مربوط به انحراف معیار $\sigma$ (شکل ۲.۷ را ببینید).

شکل ۲.۷: یک توزیع نرمال تک‌متغیره تقریباً $95\%$ از مساحت خود را در محدوده $|x-\mu| \le 2\sigma$ دارد، همانطور که نشان داده شده است. اوج توزیع دارای مقدار $p(\mu) = \frac{1}{\sqrt{2\pi}\sigma}$ است.

\section*{چگالی نرمال چندمتغیره (\lr{Multivariate Normal Density})}
بردار میانگین ($\mu$) و ماتریس کوواریانس ($\Sigma$) توزیع را توصیف می‌کنند:
\[
p(x) = \frac{1}{(2\pi)^{d/2}|\Sigma|^{1/2}} \exp\left[ -\frac{1}{2}(x-\mu)^T\Sigma^{-1}(x-\mu) \right]
\]

\subsection*{تعاریف رسمی}
به طور رسمی، داریم:
\[
\mu = E[x] = \int x p(x) dx
\]
و
\[
\Sigma = E[(x-\mu)(x-\mu)^T]
\]
که در آن مقدار مورد انتظار یک بردار یا یک ماتریس با گرفتن مقادیر مورد انتظار مولفه‌های آن به دست می‌آید. به عبارت دیگر، اگر $x_i$ مولفه $i$-ام $x$ باشد، $\mu_i$ مولفه $i$-ام $\mu$ و $\sigma_{ij}$ مولفه $ij$-ام $\Sigma$ باشد، آنگاه:
\[
\mu_i = E[x_i]
\]
و
\[
\sigma_{ij} = E[(x_i - \mu_i)(x_j - \mu_j)]
\]

\section*{ویژگی‌های ماتریس کوواریانس (\lr{Covariance Matrix})}
\begin{itemize}
    \item ماتریس کوواریانس همیشه متقارن و نیمه معین مثبت است.
    \item ما توجه خود را به حالتی محدود می‌کنیم که $\Sigma$ معین مثبت باشد، به طوری که دترمینان $\Sigma$ کاملاً مثبت باشد.
    \item عناصر روی قطر $\sigma_{ii}$ واریانس‌های $x_i$ مربوطه هستند (یعنی $\sigma_i^2$).
    \item عناصر خارج از قطر $\sigma_{ij}$ کوواریانس‌های $x_i$ و $x_j$ هستند.
\end{itemize}

مثال: برای ویژگی‌های طول و وزن یک جمعیت از ماهی‌ها، انتظار یک کوواریانس مثبت را داریم.

\begin{itemize}
    \item اگر $x_i$ و $x_j$ به صورت آماری مستقل باشند، آنگاه $\sigma_{ij} = 0$.
    \item اگر تمام عناصر خارج از قطر صفر باشند، $\Sigma$ به حاصل‌ضرب چگالی‌های نرمال تک‌متغیره برای مولفه‌های $x$ کاهش می‌یابد.
\end{itemize}

چگالی نرمال چندمتغیره به طور کامل با $d + \frac{d(d+1)}{2}$ پارامتر مشخص می‌شود:
\begin{itemize}
    \item $d$ عنصر بردار میانگین $\mu$.
    \item $\frac{d(d+1)}{2}$ عنصر مستقل ماتریس کوواریانس $\Sigma$.
\end{itemize}

نمونه‌های گرفته شده از یک جمعیت نرمال تمایل دارند در یک خوشه واحد قرار گیرند (شکل ۲.۹ را ببینید):
\begin{itemize}
    \item مرکز خوشه توسط بردار میانگین $\mu$ تعیین می‌شود.
    \item شکل خوشه توسط ماتریس کوواریانس $\Sigma$ تعیین می‌شود.
    \item مکان هندسی نقاط با چگالی ثابت، بیضی‌گون‌های با ابعاد بالا (\lr{hyperellipsoids}) هستند که توسط فرم درجه دوم زیر تعریف می‌شوند:
    \[
    (x-\mu)^T\Sigma^{-1}(x-\mu) = \text{constant}
    \]
    \item محورهای اصلی این بیضی‌گون‌های با ابعاد بالا توسط بردارهای ویژه $\Sigma$ (که با $\Phi$ نشان داده می‌شوند) داده می‌شوند.
    \item مقادیر ویژه (که با $\Lambda$ نشان داده می‌شوند) طول این محورها را تعیین می‌کنند.
    \item مقدار $r^2 = (x-\mu)^T\Sigma^{-1}(x-\mu)$ فاصله ماهالانوبیس مربع (\lr{squared Mahalanobis distance}) از $x$ تا $\mu$ نامیده می‌شود. بنابراین:
    \[
    r^2 = (x-\mu)^T\Sigma^{-1}(x-\mu)
    \]
    \item خطوط تراز چگالی ثابت، بیضی‌گون‌های با ابعاد بالا با فاصله ماهالانوبیس ثابت تا $\mu$ هستند.
    \item حجم این بیضی‌گون‌های با ابعاد بالا، پراکندگی نمونه‌ها در اطراف میانگین را اندازه‌گیری می‌کند.
\end{itemize}
\begin{figure}[h!]
  \centering
  \includegraphics[width=0.75\textwidth]{img_page_10_1.png}
  \caption{توزیع نرمال دوبعدی و بردارهای میانگین و ماتریس کوواریانس}
\end{figure}




شکل ۲.۹: نمونه‌های گرفته شده از یک توزیع گاوسی دو بعدی در ابری به مرکز بردار میانگین $\mu$ قرار دارند. بیضی‌های قرمز خطوط چگالی احتمال برابر از توزیع گاوسی را نشان می‌دهند.

\section*{توابع تشخیص‌دهنده برای چگالی نرمال}

\subsection*{حالت ۱: ماتریس‌های کوواریانس برابر ($\Sigma_i = \sigma^2 I$)}
تابع تشخیص‌دهنده خطی است:
\[
g_i(x) = w_i^T x + w_{i0}
\]

\subsubsection*{چگونه؟}
بیایید شکل ساده‌شده درست‌نمایی لگاریتمی $g_i(x)$ را برای حالتی که ماتریس‌های کوواریانس برابر و همسانگرد هستند ($\Sigma_i = \sigma^2 I$) استخراج کنیم. ما با چگالی نرمال چندمتغیره داده شده شروع می‌کنیم و نشان می‌دهیم که چگونه به فرم بیان شده کاهش می‌یابد.

تابع چگالی احتمال (\lr{PDF}) برای کلاس $\omega_i$ به صورت زیر است:
\[
p(x|\omega_i) = \frac{1}{(2\pi)^{d/2}|\Sigma_i|^{1/2}} \exp\left[ -\frac{1}{2}(x-\mu_i)^T\Sigma_i^{-1}(x-\mu_i) \right]
\]
که در آن:
\begin{itemize}
    \item $\Sigma_i = \sigma^2 I$ (کوواریانس همسانگرد، برابر در سراسر کلاس‌ها).
    \item $|\Sigma_i| = (\sigma^2)^d$ (دترمینان $\Sigma_i$).
    \item $\Sigma_i^{-1} = \frac{1}{\sigma^2} I$ (معکوس ماتریس قطری).
\end{itemize}

با جایگذاری $\Sigma_i = \sigma^2 I$ در \lr{PDF}:
\[
p(x|\omega_i) = \frac{1}{(2\pi)^{d/2}(\sigma^2)^{d/2}} \exp\left[ -\frac{1}{2}(x-\mu_i)^T \left( \frac{1}{\sigma^2} I \right)(x-\mu_i) \right]
\]

ساده‌سازی نما:
\[
(x-\mu_i)^T \left( \frac{1}{\sigma^2} I \right)(x-\mu_i) = \frac{1}{\sigma^2} (x-\mu_i)^T(x-\mu_i) = \frac{\|x-\mu_i\|^2}{\sigma^2}
\]
که در آن $\|x-\mu_i\|^2$ فاصله اقلیدسی مربع است.

بنابراین، \lr{PDF} به صورت زیر در می‌آید:
\[
p(x|\omega_i) = \frac{1}{(2\pi\sigma^2)^{d/2}} \exp\left[ -\frac{\|x-\mu_i\|^2}{2\sigma^2} \right]
\]

اعمال قاعده بیز برای احتمال پسین (\lr{Posterior Probability}):
احتمال پسین $P(\omega_i|x)$ به صورت زیر است:
\[
P(\omega_i|x) = \frac{p(x|\omega_i)P(\omega_i)}{p(x)}
\]
که در آن $p(x) = \sum_j p(x|\omega_j)P(\omega_j)$ شواهد (\lr{evidence}) است (که در توابع تشخیص‌دهنده نادیده گرفته می‌شود).

لگاریتم پسین (تابع تشخیص‌دهنده $g_i(x)$) به صورت زیر است:
\[
g_i(x) = \ln p(x|\omega_i) + \ln P(\omega_i)
\]

گرفتن لگاریتم از \lr{PDF} و محاسبه $\ln p(x|\omega_i)$:
\[
\ln p(x|\omega_i) = \ln\left( \frac{1}{(2\pi\sigma^2)^{d/2}} \right) - \frac{\|x-\mu_i\|^2}{2\sigma^2}
\]

ساده‌سازی جمله اول:
\[
\ln((2\pi\sigma^2)^{-d/2}) = -\frac{d}{2} \ln(2\pi\sigma^2)
\]
بنابراین:
\[
\ln p(x|\omega_i) = -\frac{d}{2} \ln(2\pi\sigma^2) - \frac{\|x-\mu_i\|^2}{2\sigma^2}
\]

ترکیب با پیشین و نادیده گرفتن مقادیر ثابت:
تابع تشخیص‌دهنده به صورت زیر در می‌آید:
\[
g_i(x) = \ln p(x|\omega_i) + \ln P(\omega_i) = -\frac{\|x-\mu_i\|^2}{2\sigma^2} - \frac{d}{2} \ln(2\pi\sigma^2) + \ln P(\omega_i)
\]

از آنجا که $-\frac{d}{2} \ln(2\pi\sigma^2)$ برای همه کلاس‌های $\omega_i$ ثابت است، تاثیری در تصمیم دسته‌بندی ندارد (ما $g_i(x)$ را در سراسر $i$ مقایسه می‌کنیم و جمله ثابت حذف می‌شود). بنابراین، ما آن را حذف کرده و می‌نویسیم:
\[
g_i(x) = -\frac{\|x-\mu_i\|^2}{2\sigma^2} + \ln P(\omega_i)
\]

مرزهای تصمیم‌گیری ابرصفحه (\lr{hyperplanes}) هستند.

اگر احتمالات پیشین برابر نباشند، فاصله مربع $\|x-\mu_i\|^2$ توسط واریانس $\sigma^2$ نرمال‌سازی شده و با اضافه کردن $\ln P(\omega_i)$ جبران می‌شود. این بدان معنی است که اگر $x$ به طور مساوی به دو بردار میانگین متفاوت نزدیک باشد، تصمیم بهینه از طبقه‌ای که احتمال پیشین بیشتری دارد حمایت خواهد کرد.

نیازی به محاسبه صریح فاصله‌ها نیست. بسط فرم درجه دوم $(x-\mu_i)^T(x-\mu_i)$ به دست می‌دهد:
\[
g_i(x) = -\frac{1}{2\sigma^2} [x^T x - 2\mu_i^T x + \mu_i^T \mu_i] + \ln P(\omega_i)
\]

جمله درجه دوم $x^T x$ برای همه $i$ها یکسان است، بنابراین می‌توان آن را به عنوان یک ثابت جمعی نادیده گرفت. این کار تابع تشخیص‌دهنده را به صورت زیر ساده می‌کند:
\[
g_i(x) = w_i^T x + w_{i0}
\]
که در آن:
\[
w_i = \frac{\mu_i}{\sigma^2}
\]
و
\[
w_{i0} = -\frac{\mu_i^T \mu_i}{2\sigma^2} + \ln P(\omega_i)
\]
در اینجا، $w_{i0}$ آستانه یا بایاس (\lr{threshold or bias}) در جهت $i$-ام نامیده می‌شود.

دسته‌بندی‌کننده‌ای که از توابع تشخیص‌دهنده خطی استفاده می‌کند، یک ماشین خطی (\lr{linear machine}) نامیده می‌شود. سطوح تصمیم‌گیری برای یک ماشین خطی، بخش‌هایی از ابرصفحه‌ها هستند که توسط معادلات خطی $g_i(x) = g_j(x)$ برای دو طبقه با بیشترین احتمال پسین تعریف می‌شوند.

برای مورد ما، این معادله را می‌توان به صورت زیر نوشت:
\[
w^T(x - x_0) = 0
\]
که در آن:
\[
w = \mu_i - \mu_j
\]
و
\[
x_0 = \frac{1}{2}(\mu_i + \mu_j) - \frac{\sigma^2}{\|\mu_i - \mu_j\|^2} \ln \frac{P(\omega_i)}{P(\omega_j)} (\mu_i - \mu_j)
\]

این معادله یک ابرصفحه را تعریف می‌کند که از نقطه $x_0$ عبور می‌کند و بر بردار $w$ عمود است. از آنجا که $w = \mu_i - \mu_j$، ابرصفحه جداکننده $R_i$ و $R_j$ بر خط اتصال میانگین‌ها عمود است. اگر $P(\omega_i) = P(\omega_j)$ باشد، ابرصفحه عمودمنصف خط بین میانگین‌ها است. اگر احتمالات پیشین نابرابر باشند، نقطه $x_0$ از میانگین محتمل‌تر دور می‌شود.

شکل ۲.۱۰: اگر کوواریانس‌های دو توزیع برابر و متناسب با ماتریس همانی باشند، توزیع‌ها در $d$ بعد کروی هستند و مرز یک ابرصفحه تعمیم‌یافته از $d-1$ بعد است که عمود بر خط جداکننده میانگین‌ها قرار دارد.

معادله $w^T(x - x_0) = 0$ نشان‌دهنده یک ابرصفحه در فضای $d$-بعدی است. نحوه تفسیر و تجسم آن به صورت زیر است:

\subsubsection*{تفسیر}
\begin{enumerate}
    \item $w$: یک بردار نرمال برای ابرصفحه (جهت ابرصفحه را تعریف می‌کند).
    \item $x_0$: یک نقطه ثابت روی ابرصفحه.
    \item $x$: یک نقطه متغیر روی ابرصفحه.
\end{enumerate}
معادله بیان می‌کند که بردار $x - x_0$ بر $w$ عمود است، به این معنی که تمام نقاط $x$ روی ابرصفحه این شرط را برآورده می‌کنند.

\subsubsection*{تجسم در دو بعد (2D)}
در فضای دو بعدی ($d=2$)، معادله به یک خط کاهش می‌یابد:
\begin{enumerate}
    \item $w = [w_1, w_2]^T$: بردار نرمال خط.
    \item $x_0 = [x_0, y_0]^T$: یک نقطه ثابت روی خط.
    \item $x = [x, y]^T$: یک نقطه متغیر روی خط.
\end{enumerate}
معادله به صورت زیر در می‌آید:
\[
w_1(x - x_0) + w_2(y - y_0) = 0
\]
این معادله یک خط در دو بعد است.

\subsubsection*{رسم ابرصفحه (خط در 2D)}
نحوه رسم آن به این صورت است:
\begin{enumerate}
    \item نقطه $x_0$ را رسم کنید: این یک نقطه ثابت روی خط است.
    \item بردار نرمال $w$ را رسم کنید: این بردار بر خط عمود است.
    \item خط را رسم کنید: خط از $x_0$ عبور می‌کند و بر $w$ عمود است.
\end{enumerate}

\subsubsection*{مثال}
بیایید از مقادیر زیر استفاده کنیم:
\begin{itemize}
    \item $w = [2, 1]^T$ (بردار نرمال).
    \item $x_0 = [1, 1]^T$ (یک نقطه روی خط).
\end{itemize}
معادله به صورت زیر در می‌آید: $2(x-1) + 1(y-1) = 0$. ساده‌سازی کنید:
\[
2x-2+y-1=0 \implies 2x+y-3=0
\]

\subsection*{ابرصفحه و مرز تصمیم‌گیری (\lr{Hyperplane and Decision Boundary})}
معادله $w^T(x-x_0) = 0$ یک ابرصفحه را تعریف می‌کند که از نقطه $x_0$ عبور می‌کند و بر بردار $w$ عمود است. از آنجا که $w = \mu_i - \mu_j$، ابرصفحه جداکننده $R_i$ و $R_j$ بر خط اتصال میانگین‌ها عمود است.

اگر $P(\omega_i) = P(\omega_j)$ باشد، جمله دوم ناپدید می‌شود و نقطه $x_0$ در نیمه راه بین میانگین‌ها قرار دارد. در این حالت، ابرصفحه عمودمنصف خط بین میانگین‌ها است (شکل ۲.۱۱ را ببینید).

اگر احتمالات پیشین نابرابر باشند، نقطه $x_0$ از میانگین محتمل‌تر دور می‌شود.

توجه داشته باشید که اگر واریانس $\sigma^2$ نسبت به فاصله مربع $\|\mu_i - \mu_j\|^2$ کوچک باشد، موقعیت مرز تصمیم‌گیری نسبت به مقادیر دقیق احتمالات پیشین نسبتاً غیرحساس است.

\subsection*{دسته‌بندی‌کننده با حداقل فاصله (\lr{Minimum Distance Classifier})}
اگر احتمالات پیشین $P(\omega_i)$ برای همه $c$ کلاس یکسان باشد، جمله $\ln P(\omega_i)$ به یک ثابت جمعی بی‌اهمیت تبدیل می‌شود و می‌توان از آن صرف‌نظر کرد.

در این حالت، قاعده تصمیم‌گیری بهینه به صورت زیر ساده می‌شود:
\begin{itemize}
    \item فاصله اقلیدسی $\|x - \mu_i\|$ را از $x$ تا هر یک از $c$ بردار میانگین اندازه‌گیری کنید.
    \item $x$ را به نزدیک‌ترین میانگین اختصاص دهید.
\end{itemize}

چنین دسته‌بندی‌کننده‌ای، دسته‌بندی‌کننده با حداقل فاصله (\lr{minimum distance classifier}) نامیده می‌شود.

اگر هر بردار میانگین به عنوان یک نمونه اولیه یا الگوی ایده‌آل برای الگوهای کلاس خود در نظر گرفته شود، این اساساً یک روش تطبیق الگو (\lr{template-matching}) است (شکل ۲.۱۰ را ببینید). این تکنیک شبیه به الگوریتم نزدیک‌ترین همسایه (\lr{nearest-neighbor algorithm}) است.

شکل ۲.۱۱: با تغییر پیشین‌ها، مرز تصمیم‌گیری تغییر می‌کند؛ برای پیشین‌های بسیار متفاوت، مرز بین میانگین‌های این توزیع‌های گاوسی کروی ۱، ۲ و ۳ بعدی قرار نخواهد گرفت.

\subsection*{حالت ۲: ماتریس‌های کوواریانس برابر اما دلخواه ($\Sigma_i = \Sigma$)}
تابع تشخیص‌دهنده همچنان خطی است:
\[
g_i(x) = -\frac{1}{2}(x-\mu_i)^T \Sigma^{-1} (x-\mu_i) + \ln P(\omega_i)
\]
مرزهای تصمیم‌گیری ابرصفحه هستند، اما لزوماً بر خط بین میانگین‌ها عمود نیستند.

\subsection*{حالت ۳: ماتریس‌های کوواریانس دلخواه ($\Sigma_i \neq \Sigma_j$)}
تابع تشخیص‌دهنده درجه دوم است:
\[
g_i(x) = x^T W_i x + w_i^T x + w_{i0}
\]
که در آن:
\begin{itemize}
    \item $W_i = -\frac{1}{2}\Sigma_i^{-1}$
    \item $w_i = \Sigma_i^{-1}\mu_i$
    \item $w_{i0} = -\frac{1}{2}\mu_i^T \Sigma_i^{-1}\mu_i - \frac{1}{2}\ln|\Sigma_i| + \ln P(\omega_i)$
\end{itemize}

\section*{خلاصه}
\begin{itemize}
    \item \lr{Bayesian Decision Theory} چارچوبی را برای دسته‌بندی بهینه تحت عدم قطعیت فراهم می‌کند.
    \item مولفه‌های کلیدی:
    \begin{itemize}
        \item احتمالات پیشین (\lr{Prior probabilities}).
        \item درست‌نمایی‌ها (\lr{Likelihoods}) (چگالی‌های شرطی کلاس).
        \item احتمالات پسین (\lr{Posterior probabilities}).
    \end{itemize}
    \item قواعد تصمیم‌گیری، نرخ‌های خطا یا زیان مورد انتظار را کمینه می‌کنند.
    \item توزیع‌های نرمال به دلیل تحلیل‌پذیری ریاضی آنها معمولاً استفاده می‌شوند.
\end{itemize}
\begin{figure}[h!]
  \centering
  \includegraphics[width=0.75\textwidth]{img_page_5_1.png}
  \caption{احتمالات پسین مربوطه با در نظر گرفتن مقادیر احتمال پیشین}
\end{figure}

\begin{figure}[h!]
  \centering
  \includegraphics[width=0.75\textwidth]{img_page_18_1.png}
  \caption{مرز تصمیم خطی (ابرصفحه) برای کوواریانس‌های برابر}
\end{figure}






\end{document}
